\chapter{Ce que les routes enseignent}
\label{chap:ce-que-les-routes-enseignent}

Les semaines devinrent des mois.

Khaer Ghal et Kamatsu avaient désormais trouvé leur place dans la caravane.
Ils participaient à la vie du camp, aidaient lorsqu'il fallait charger ou
décharger les chariots, s'occupaient des animaux et accomplissaient les mêmes
corvées que les autres. À chaque étape, ils découvraient de nouveaux paysages,
de nouvelles villes et surtout de nouvelles personnes.

Car la caravane n'était jamais tout à fait la même.

Certains voyageurs l'accompagnaient pendant quelques jours avant de prendre
une autre route. D'autres restaient plusieurs semaines ou plusieurs mois.
Marchands, artisans, soldats, voyageurs ou simples vagabonds partageaient un
temps les mêmes feux avant de poursuivre leur chemin.

Pour les deux garçons, chacune de ces rencontres pouvait devenir une occasion
d'apprendre quelque chose.

Ils grandissaient ensemble, mais leurs curiosités commençaient déjà à prendre
des directions différentes.

Khaer Ghal observait beaucoup. Il s'intéressait volontiers à ce qui pouvait
lui être utile sur les routes : une manière différente de tenir une arme, une
trace qu'il ne connaissait pas, un nœud, une plante, une façon de se déplacer
ou simplement un geste qu'il n'avait jamais vu auparavant.

Kamatsu, lui, posait davantage de questions. Il s'intéressait aux histoires
que racontaient les voyageurs, aux endroits qu'ils avaient visités et aux
choses qu'ils savaient fabriquer. Là où Khaer Ghal pouvait passer une heure à
observer quelqu'un travailler sans prononcer un mot, Kamatsu avait généralement
posé dix questions avant même que celui-ci ait terminé.

Ils apprenaient différemment.

Mais ils apprenaient tous les deux.

\scenebreak

Khaer Ghal possédait déjà certaines connaissances acquises dans son village.
Comme beaucoup de jeunes Orcs de son clan, il avait reçu des rudiments de
combat, appris à manier une dague et une lame courte, à se défendre contre
certaines créatures et à reconnaître les moments où il fallait frapper,
esquiver ou profiter d'une ouverture.

À cela s'ajoutait désormais l'expérience beaucoup plus concrète de ce qu'il
avait réellement vécu.

Kamatsu avait beaucoup moins d'expérience dans ce domaine.

Un soir, lors d'une halte, Khaer Ghal décida donc de lui montrer ce qu'il
savait.

Les deux garçons trouvèrent des bâtons et Khaer Ghal entreprit très
sérieusement de lui expliquer comment réagir face à une créature qui attaquait.

— Non. Pas comme ça. Si elle te saute dessus, tu n'auras pas le temps.

Kamatsu baissa légèrement son bâton.

— Si quoi me saute dessus ?

— La créature.

— Quelle créature ?

Khaer Ghal réfléchit.

— Celle qui t'attaque.

Kamatsu regarda autour de lui.

— Il n'y a pas de créature.

— Mais imagine.

Khaer Ghal lui montra comment garder son arme près du corps, attendre que
l'animal approche, s'écarter lorsqu'il bondissait puis profiter de l'ouverture
pour frapper.

Kamatsu observa attentivement le mouvement.

— Mais toi, tu as un bâton.

— Imagine que c'est une dague.

— Alors pourquoi j'ai un bâton ?

— Parce qu'on n'a pas de dagues.

La réponse semblait parfaitement logique à Khaer Ghal.

Un peu moins à Kamatsu.

Ce qu'aucun des deux n'avait remarqué, c'était le soldat assis contre un
chariot quelques mètres plus loin. Il les observait depuis un moment,
Bartiméus confortablement installé sur ses genoux.

— Tu en penses quoi ? demanda-t-il au chat.

Bartiméus ne bougea pas.

— Oui. Moi aussi.

Le jeune Orc ne racontait pourtant pas n'importe quoi. Ses explications étaient
maladroites et ses mouvements parfois approximatifs, mais le soldat
reconnaissait derrière tout cela quelque chose de bien réel : l'expérience.
Khaer Ghal ne savait pas véritablement manier une épée ; il savait se défendre
contre quelque chose qui voulait le tuer, et ce n'était pas tout à fait la même
chose.

Lorsque Kamatsu reprit sa garde, le soldat intervint.

— Et si elle a une épée ?

Khaer Ghal se retourna.

— Une créature avec une épée ?

— Non. Un homme.

Khaer Ghal regarda son bâton, puis Kamatsu.

— Je sais pas.

Le soldat sourit.

— Voilà une bonne réponse.

Il déposa Bartiméus, qui protesta contre cette décision prise sans
consultation, puis ramassa un troisième bâton. Il ne leur enseigna rien de
spectaculaire : simplement les pieds, la garde, la distance, la façon
d'accompagner une parade plutôt que d'opposer toute sa force à un coup. Il
leur expliqua surtout qu'une technique efficace contre une bête ne l'était pas
nécessairement contre quelqu'un capable d'anticiper et de réfléchir.

À la fin de la leçon, il leur proposa de le toucher. Khaer Ghal attaqua le
premier et, quelques secondes plus tard, son bâton reposait par terre. Kamatsu
éclata de rire, jusqu'à ce que le soldat lui fasse signe d'essayer à son tour.
Peu après, deux bâtons étaient au sol et Khaer Ghal riait à son tour.

Le soldat récupéra Bartiméus et se rassit.

— Maintenant, recommencez.

Les garçons ramassèrent leurs armes. Cette fois, leur jeu ressemblait déjà un
peu moins à un jeu.

\scenebreak

Ce ne fut ni la dernière leçon, ni la dernière personne auprès de laquelle ils
apprirent quelque chose.

Sur les routes, les occasions ne manquaient pas.

Un chasseur pouvait montrer comment reconnaître une piste presque effacée. Un
marchand pouvait expliquer pourquoi deux pièces qui semblaient identiques
n'avaient pas la même valeur. Un artisan pouvait leur apprendre à réparer une
sangle plutôt que de la jeter. Un voyageur connaissait une plante que Khaer
Ghal n'avait encore jamais vue ; un autre savait prévoir l'arrivée de la pluie
en observant le ciel.

Même les tâches les plus ordinaires finissaient par devenir des enseignements.

Khaer Ghal ne le savait pas encore, mais les routes allaient lui offrir quelque
chose que son village n'aurait jamais pu lui donner. Dans la caravane, les gens
arrivaient et repartaient, chacun avec ses connaissances, ses habitudes et sa
manière de voir le monde. Il n'était pas obligé de choisir un seul maître ni
d'attendre que quelqu'un décide de ce qu'il devait apprendre : il pouvait
regarder, écouter, essayer, se tromper, recommencer et conserver de chaque
rencontre ce qui lui permettait d'avancer.

Pour un garçon aussi curieux que lui, les routes étaient devenues une école
immense.

Et il avait bien l'intention d'en profiter.